\chapter{Private Connectivity to Crypto Venues}\label{ch:nw:private-connectivity-to-crypto-venues}

A venue that runs in a \omterm{def:nw:trading-in-a-public-cloud:cloud}{public cloud} can offer its market makers a \omterm{def:nw:private-connectivity-to-crypto-venues:endpoint}{private endpoint} inside the provider's network: their orders never touch
the public internet, and, if the endpoint has an interface in the same zone as their machines, they never cross a zone boundary either. The
zone identifiers of chapter~16 let a firm check which is the case. The difference is not small: in this chapter's model, the same \omterm{def:nw:private-connectivity-to-crypto-venues:endpoint}{private
endpoint} served from another zone adds about \qty{285}{\micro\second} to every median round trip, and costs less than a cent more per million
messages. Private connectivity is cheap; getting it in the wrong place is expensive only in time.

Chapter~17 found the venue's region and zone; this chapter connects to it. Virtual networks and the two ways of joining them, peering and
\omterm{def:nw:private-connectivity-to-crypto-venues:endpoint}{private endpoints}; what the venues offer (from shared cluster placement to \omterm{def:nw:colocation-products-and-how-they-are-sold:colo}{colocation} \omterm{def:nw:colocation-products-and-how-they-are-sold:mmr}{cross-connects} for the few that run in data centres);
and the tiers of access that unlock the faster paths. The venues' documentation and the provider's are cited and dated; the latencies are a
labelled simulation built on chapter~16's published measurement.

\section{Virtual networks, peering and private endpoints}

\begin{definition}[Virtual private cloud, virtual-network peering]\label{def:nw:private-connectivity-to-crypto-venues:peering}
A \emph{virtual private cloud}\index{virtual private cloud} (VPC) is a customer's private network inside a \omterm{def:nw:trading-in-a-public-cloud:cloud}{public cloud}: its own address space,
subnets in chosen zones, routing and access rules. \emph{Virtual-network peering}\index{virtual-network peering} joins two such networks, of the
same or different customers, so that their machines exchange traffic over private addresses as if they were one network.
\end{definition}

\begin{definition}[Private endpoint, network load balancer]\label{def:nw:private-connectivity-to-crypto-venues:endpoint}
A \emph{private endpoint}\index{private endpoint} is a network interface that a provider places in a customer's virtual network, in chosen zones,
through which the customer reaches another party's service privately, without routing between the two networks. A \emph{network load
balancer}\index{network load balancer} is the provider's service that receives connections on one address and spreads them over target machines,
at the level of TCP or UDP; a private-endpoint service is typically exposed through one.
\end{definition}

Peering and endpoints solve different problems. Peering joins two networks and lets them route to each other: fast, simple, but both sides must
agree on addresses and trust each other's routes. An endpoint publishes one service to many customers without joining their networks: each
customer gets an interface in its own network, and the provider carries its packets through a load balancer to the service. AWS documents
both: peering makes instances ``communicate with each other as if they are within the same network''; PrivateLink connects a network to a
service ``as if they were in your VPC''.

\begin{omfigure}
\begin{tikzpicture}[font=\footnotesize,b/.style={draw,rounded corners=2pt,minimum height=0.55cm,align=center,inner sep=3pt},>=Stealth]
\node[b,fill=omDef!15,minimum width=2.4cm] (f) at (0,0) {firm's instance\\ (zone az4)};
\node[b,fill=omThm!15,minimum width=2.6cm] (v) at (10.4,0) {venue's gateways\\ (its VPC)};
\node[b,fill=gray!10] (e) at (3.4,-1.5) {private endpoint\\ interface};
\node[b,fill=gray!10] (l) at (7,-1.5) {network load\\ balancer};
\draw[->,thick] (f) -- node[above,font=\scriptsize]{peering: route between the two VPCs} (v);
\draw[->,thick,omMeth] (f) |- (e); \draw[->,thick,omMeth] (e) -- (l); \draw[->,thick,omMeth] (l) -| (v);
\node[font=\scriptsize,omMeth] at (5.2,-2.25) {private endpoint: extra hops, no routing between the networks};
\node[b,fill=omMeth!10] (x) at (5.2,1.6) {edge network (public)};
\draw[->,thick,gray,dashed] (f) |- (x); \draw[->,thick,gray,dashed] (x) -| (v);
\end{tikzpicture}

\omcaption{Three ways from a firm's instance to a venue in the same cloud region, schematically: peering between the two virtual networks, a \omterm{def:nw:private-connectivity-to-crypto-venues:endpoint}{private
endpoint} (an interface in the firm's network, a load balancer in front of the venue's gateways), and the public endpoint through an \omterm{def:nw:where-the-crypto-venues-live:as}{edge network}.}\label{fig:nw:private-connectivity-to-crypto-venues:paths}
\end{omfigure}

The endpoint's zone is the catch. An endpoint has interfaces in the zones where it is created; the load
balancer behind it sends each connection to targets in its zone unless cross-zone balancing is enabled, which AWS documents as off by default
for \omterm{def:nw:private-connectivity-to-crypto-venues:endpoint}{network load balancers}. If the venue's gateways are in one physical zone and the firm's instance in another, every packet crosses the
boundary somewhere; the only question is where. Chapter~16's identifiers let the firm compare its instance's zone with the endpoint's
interfaces before it sends a single order.

\begin{proposition}[Where the zone boundary lies]\label{prop:nw:private-connectivity-to-crypto-venues:zone}
If the firm's instance and the venue's gateways are in different physical zones, every path between them crosses a zone boundary, whatever the
endpoint's placement; if they are in the same zone, a path crosses one only if the endpoint or the load balancer routes the traffic through
another zone. The firm therefore needs both identifiers (its instance's and the venue's gateways') to know whether a same-zone path exists.
\end{proposition}

\begin{proof}
A path joins two points; if they lie in different zones some link of the path joins the zones. If they lie in one zone, only a hop placed in
another zone (an endpoint interface or a load-balancer node) takes the path out of it and back.
\end{proof}

\begin{omfigure}
\begin{tikzpicture}[font=\footnotesize,b/.style={draw,rounded corners=2pt,minimum height=0.5cm,align=center,inner sep=2pt},>=Stealth]
\foreach \x/\z in {0/az1,4.2/az2,8.4/az4}{\draw[gray,rounded corners=3pt] (\x,-1.3) rectangle (\x+3.6,1.5); \node[gray,font=\small] at (\x+1.8,1.25) {\z};}
\node[b,fill=omThm!20] (g) at (1.8,0.4) {venue's gateways};
\node[b,fill=gray!10] (e1) at (1.8,-0.6) {endpoint interface};
\node[b,fill=gray!10] (e2) at (6.0,-0.6) {endpoint interface};
\node[b,fill=omDef!15] (i) at (10.2,0.4) {firm's instance};
\draw[->,thick,omMeth] (i) .. controls (8,-0.2) and (7.5,-0.6) .. (e2);
\draw[->,thick,omMeth] (e2.west) .. controls (4.0,-0.6) and (3.6,0.4) .. (g.east);
\node[font=\scriptsize,omMeth,align=center] at (6.0,0.6) {every packet crosses\\ two zone boundaries};
\end{tikzpicture}

\omcaption{The endpoint in the wrong zone, schematically: the venue's gateways in az1, its endpoint's interfaces in az1 and az2, the firm's instance in
az4; the firm's packets enter the interface in az2 and cross to the gateways through the load balancer. No placement of the endpoint keeps this path in one zone (\cref{prop:nw:private-connectivity-to-crypto-venues:zone}); moving the instance to az1
does. Zone labels are identifiers, not names.}\label{fig:nw:private-connectivity-to-crypto-venues:zones}
\end{omfigure}

\omcode{code/firm/privlink/firm_privlink.py}{59}{77}{A path's round trip (the body of chapter~16's model plus the documented extra hops), the zone
check, and the costs.}\label{lst:nw:private-connectivity-to-crypto-venues:path}

\begin{dated}{2026-09}{Private-endpoint prices}\label{dat:nw:private-connectivity-to-crypto-venues:prices}
AWS's price list of 17 September 2026 charges 0.014~USD an hour for each interface endpoint in each zone in Asia Pacific (Tokyo), 0.01 in US East
(N. Virginia), and 0.01~USD per gigabyte of data processed for the first petabyte a month in a region (0.006 for the next four, 0.004 beyond); data
crossing zones is charged 0.01~USD per gigabyte in each direction (chapter~16).
\end{dated}

\section{Venue colocation offers}

\begin{dated}{2026-09}{What venues offer, from their documentation}\label{dat:nw:private-connectivity-to-crypto-venues:offers}
\begin{itemize}
\item Coinbase Derivatives: production in Equinix CH4 (350 E Cermak, Chicago), disaster recovery and integration in Equinix NY5 (Secaucus); access by
  \omterm{def:nw:colocation-products-and-how-they-are-sold:colo}{colocation} \omterm{def:nw:colocation-products-and-how-they-are-sold:mmr}{cross-connect} (recommended for the lowest and most stable latency), by AWS PrivateLink (``some of the benefits of private connectivity
  without having to cross connect''), or over the internet with TLS.
\item One Trading, with AWS (December 2023): four access tiers tested for market makers, from VPC peering inside a cluster \omterm{def:nw:trading-in-a-public-cloud:pg}{placement group} shared
  between the exchange's and the market maker's accounts (lowest latency), through plain peering and PrivateLink behind a \omterm{def:nw:private-connectivity-to-crypto-venues:endpoint}{network load balancer}
  (medium, with extra hops), to the internet through an \omterm{def:nw:where-the-crypto-venues-live:as}{edge network}; the exchange's stated goal, round trips under \qty{200}{\micro\second}.
\end{itemize}
\end{dated}

A shared cluster \omterm{def:nw:trading-in-a-public-cloud:pg}{placement group} is the cloud's version of \omterm{def:nw:colocation-products-and-how-they-are-sold:colo}{colocation}: the venue creates the group in its account and shares it, so that the
market maker's instances land in the same high-bisection segment of the network as the venue's gateways and matching engine. It is the only
offer on the list that brings the customer's machine physically closer; the others choose a route to wherever the machine happens to be.

\section{Cross-connects at data-centre-hosted venues}

Some crypto venues do not run in a cloud at all. Coinbase's derivatives exchange runs in a Chicago data centre with its disaster recovery in New
Jersey, the same buildings as chapter~10's venues, and sells the connections of chapter~9: \omterm{def:nw:colocation-products-and-how-they-are-sold:mmr}{cross-connects} to the exchange's network in the building.
Its documentation is explicit about the hierarchy: \omterm{def:nw:colocation-products-and-how-they-are-sold:mmr}{cross-connects} for the lowest latency, PrivateLink for private connectivity without a
\omterm{def:nw:colocation-products-and-how-they-are-sold:mmr}{cross-connect}, the internet for the rest. For such a venue the cloud is an access network, not the venue's home, and the firm's choice is chapter~9's.

\section{Dedicated and market-maker gateways, and the tiers that unlock them}

The faster paths are not open to everyone. Venues reserve shared \omterm{def:nw:trading-in-a-public-cloud:pg}{placement groups}, \omterm{def:nw:private-connectivity-to-crypto-venues:endpoint}{private endpoints} and dedicated gateways for customers above a
volume or a fee, as One Quant Book~3's dated box on dedicated endpoints and tier matching describes, and may apply separate rate limits to them.
For the firm that qualifies, the question is the one of this chapter: which path, from which zone. For the one that does not, it is chapter~17's:
the best instance it can find on the public path.

\begin{table}[ht]
\centering\small
\begin{tabular}{lrrrr}
\toprule
Path & p50 & p99 & Cost per million & Fixed per zone \\
 & (\si{\micro\second}) & (\si{\micro\second}) & messages (USD) & (USD a month) \\
\midrule
Peering, shared cluster placement & 150 & 220 & 0 & 0 \\
Peering, same zone & 270 & 396 & 0 & 0 \\
\omterm{def:nw:private-connectivity-to-crypto-venues:endpoint}{Private endpoint}, same zone & 310 & 436 & 0.004 & 10.22 \\
\omterm{def:nw:private-connectivity-to-crypto-venues:endpoint}{Private endpoint}, other zone & 595 & 797 & 0.012 & 10.22 \\
Public endpoint via an edge & 2\,270 & 2\,396 & 0 & 0 \\
\bottomrule
\end{tabular}
\caption{Simulation of the documented paths in one region: chapter~16's fitted round trips (the shared cluster taken at the AWS blog's ``sub-150~µs''),
plus an assumed \qty{40}{\micro\second} for the endpoint and load-balancer hops and \qty{2}{\milli\second} for the public edge; costs from the dated
prices, for 400 bytes per message (an order and its acknowledgement). Data: \texttt{nw\_private.path\_rows()}.}\label{tab:nw:private-connectivity-to-crypto-venues:paths}
\end{table}

\begin{omfigure}
\begin{tikzpicture}
\begin{axis}[ybar,width=11.5cm,height=6cm,bar width=9pt,ymode=log,ymin=100,ymax=5000,log origin=infty,xtick=data,
  symbolic x coords={peering + shared cluster,peering,endpoint same zone,endpoint other zone,public via edge},
  x tick label style={font=\scriptsize,text width=1.8cm,align=center},ylabel={round trip (\si{\micro\second}, log)},
  label style={font=\small},grid=major,grid style={gray!25},legend style={font=\footnotesize},legend pos=north west,
  table/col sep=comma,enlarge x limits=0.12]
\addplot[fill=omDef!55,draw=omDef] table[x=path,y=p50]{figdata/networks/18-private-connectivity-to-crypto-venues/paths.csv};
\addplot[fill=omThm!45,draw=omThm] table[x=path,y=p99]{figdata/networks/18-private-connectivity-to-crypto-venues/paths.csv};
\legend{p50,p99}
\end{axis}
\end{tikzpicture}

\omcaption{Simulation: median and 99th-percentile round trips of the five paths of \cref{tab:nw:private-connectivity-to-crypto-venues:paths}. The endpoint
in the wrong zone costs as much as the step from peering to the endpoint several times over; the public edge costs milliseconds. Data:
\texttt{fig\_private.py}.}\label{fig:nw:private-connectivity-to-crypto-venues:rtt}
\end{omfigure}

\begin{omfigure}
\begin{tikzpicture}
\begin{axis}[width=11cm,height=5.4cm,xmin=0,xmax=10000,ymin=0,ymax=140,xlabel={messages a month (millions)},x tick label style={/pgf/number format/1000 sep={\,}},
  ylabel={endpoint cost (USD a month)},tick label style={font=\footnotesize},label style={font=\small},grid=major,grid style={gray!25},
  legend pos=north west,legend style={font=\footnotesize},table/col sep=comma,scaled x ticks=false]
\addplot[omDef,thick,no markers] table[x=millions,y=same]{figdata/networks/18-private-connectivity-to-crypto-venues/monthly.csv};
\addplot[omThm,thick,dashed,no markers] table[x=millions,y=other]{figdata/networks/18-private-connectivity-to-crypto-venues/monthly.csv};
\legend{endpoint in the same zone,endpoint in another zone}
\end{axis}
\end{tikzpicture}

\omcaption{The monthly cost of one \omterm{def:nw:private-connectivity-to-crypto-venues:endpoint}{private endpoint} in one zone against message volume, from the dated prices: 10.22~USD of endpoint-hours plus data
processing, and cross-zone transfer when the endpoint serves another zone. Ten billion messages a month cost 50 or 130~USD. Data:
\texttt{fig\_private.py}.}\label{fig:nw:private-connectivity-to-crypto-venues:cost}
\end{omfigure}

\omcode{code/networks/18-private-connectivity-to-crypto-venues/python/nw_private.py}{32}{36}{The named result: what the other zone adds to the
round trip and to the cost per million messages.}\label{lst:nw:private-connectivity-to-crypto-venues:wrong}

The costs are an afterthought: one endpoint in one zone for a month is 10.22~USD of endpoint-hours, and even ten billion messages add 40~USD of data
processing, 120 across zones. What matters is the zone. The endpoint in the wrong zone adds \qty{285}{\micro\second} to the median round trip and
\qty{360}{\micro\second} to the 99th percentile, more than the whole difference between peering and a same-zone endpoint.

\begin{method}[Connecting privately to a cloud-hosted venue]\label{met:nw:private-connectivity-to-crypto-venues:connect}
\begin{enumerate}
\item Read the venue's connectivity documentation: which paths it offers (shared placement, peering, endpoints, \omterm{def:nw:colocation-products-and-how-they-are-sold:mmr}{cross-connects}), to whom, in which
  zones, at what rate limits; record each with its date.
\item Resolve the zone identifiers of your instances and of the venue's endpoint interfaces; place instances in the zone of the venue's gateways.
\item Prefer the path with the fewest hops the venue allows; check that the endpoint has an interface in your zone and that its load balancer keeps
  traffic in the zone.
\item Measure each path's round trip with hardware timestamps; keep the public path as a fallback, and know its latency.
\item Price the endpoint-hours and data processing, and check them against the latency; they will be small.
\end{enumerate}
\end{method}

\section{Tutorial: which path, which zone}\label{tut:nw:private-connectivity-to-crypto-venues:tut}

\begin{tutorial}
\textbf{Goal.} Compare the documented paths to a cloud-hosted venue by round trip and cost, and check zone alignment.
\textbf{End state:} \cref{fig:nw:private-connectivity-to-crypto-venues:rtt,fig:nw:private-connectivity-to-crypto-venues:cost} and
\cref{tab:nw:private-connectivity-to-crypto-venues:paths}.
\begin{tutsteps}
\item \textbf{Paths.} \texttt{firm\_privlink.load\_paths} reads \texttt{data/paths.csv}: each documented path with its extra hops, its prices and its
  sources.
\item \textbf{Round trips.} \texttt{rtt\_percentiles} (\cref{lst:nw:private-connectivity-to-crypto-venues:path}) adds the hops to chapter~16's body
  models.
\item \textbf{Alignment.} \texttt{aligned} compares the instance's zone identifier with the endpoint's; \texttt{nw\_private.ALIGN} is a failing example.
\item \textbf{Costs.} \texttt{cost\_per\_million} and \texttt{monthly\_fixed}; \texttt{monthly\_curve} draws the cost against volume.
\end{tutsteps}
\textbf{What to change next.} Replace the assumed \qty{40}{\micro\second} of endpoint hops with a measurement; enable cross-zone load balancing in the
model and route a fraction of connections to the other zone.
\end{tutorial}

\section{Build: the private-path planner}\label{bld:nw:private-connectivity-to-crypto-venues:privlink}

\begin{build}
\bfield{Purpose} The paths from the firm's cloud machines to each venue, their round trips and costs, and whether they stay in one zone: the input of
the cloud rows of chapter~29's plan.
\bfield{Interface} \texttt{firm\_privlink}: \texttt{Path}, \texttt{load\_paths}, \texttt{BODIES}, \texttt{rtt\_percentiles}, \texttt{aligned},
\texttt{cost\_per\_million}, \texttt{monthly\_fixed}; on \texttt{firm.cloudplan}.
\bfield{Rules} A path exists only if a venue's documentation offers it, with its source and date; extra hops are stated assumptions until measured; an
unaligned endpoint is reported, never silently accepted.
\bfield{Acceptance tests} \texttt{code/firm/privlink/tests/}: the catalogue's sources, the alignment check, the order of the paths' medians, and the
costs by hand.
\bfield{Stretch} Read endpoint interfaces and their zones from the provider's interface; add the venue's rate limits per path.
\end{build}

\begin{omsources}
\item AWS documentation: VPC peering, PrivateLink and its pricing, Network Load Balancers; AWS public price list (September 2026).
\item Coinbase Developer Documentation, Derivatives connectivity; AWS for Industries blog, ``One Trading and AWS: Cloud-native colocation for crypto
trading'' (December 2023).
\end{omsources}

\section{Exercises}

\begin{exercise}[$\star$]\label{exo:nw:private-connectivity-to-crypto-venues:1}
What does one \omterm{def:nw:private-connectivity-to-crypto-venues:endpoint}{private endpoint} cost a month in one zone in Tokyo, before any data?
\end{exercise}

\begin{exercise}[$\star$]\label{exo:nw:private-connectivity-to-crypto-venues:2}
What is the difference between peering and a \omterm{def:nw:private-connectivity-to-crypto-venues:endpoint}{private endpoint}, for the venue?
\end{exercise}

\begin{exercise}[$\star$]\label{exo:nw:private-connectivity-to-crypto-venues:3}
Why does a shared cluster \omterm{def:nw:trading-in-a-public-cloud:pg}{placement group} bring the market maker closer, while an endpoint does not?
\end{exercise}

\begin{exercise}[$\star\star$]\label{exo:nw:private-connectivity-to-crypto-venues:4}
Your instance is in apne1-az4; the venue's endpoint has interfaces in apne1-az1 and apne1-az2. What do you do?
\end{exercise}

\begin{exercise}[$\star\star$]\label{exo:nw:private-connectivity-to-crypto-venues:5}
Why is the public path through an \omterm{def:nw:where-the-crypto-venues-live:as}{edge network} slower than every private one, even inside the same region?
\end{exercise}

\begin{exercise}[$\star\star$]\label{exo:nw:private-connectivity-to-crypto-venues:6}
A venue runs in a Chicago data centre and offers PrivateLink. Why would a firm in the same data centre still buy a \omterm{def:nw:colocation-products-and-how-they-are-sold:mmr}{cross-connect}?
\end{exercise}

\begin{exercise}[$\star\star\star$]\label{exo:nw:private-connectivity-to-crypto-venues:7}
\emph{Coding.} With \texttt{firm.privlink}, find the endpoint-hop time at which a same-zone endpoint's median equals plain peering's 99th percentile.
\end{exercise}

\begin{exercise}[$\star\star\star$]\label{exo:nw:private-connectivity-to-crypto-venues:8}
\emph{Find the flaw.} ``We use the venue's \omterm{def:nw:private-connectivity-to-crypto-venues:endpoint}{private endpoint}, so our orders take the fastest path available.''
\end{exercise}

\section{Problem: The Endpoint in the Wrong Zone}

\begin{problem}[{Weekend problem --- what the wrong zone costs}]\label{pb:nw:private-connectivity-to-crypto-venues:1}
A market maker reaches a venue through the venue's \omterm{def:nw:private-connectivity-to-crypto-venues:endpoint}{private endpoint}. Its instances are in one zone; the endpoint serves it from another. Messages are
400 bytes (an order and its acknowledgement); the model's endpoint hops take \qty{40}{\micro\second}.

\medskip\noindent\textbf{Part I --- The round trip.}
\begin{enumerate}
\item What are the median and 99th-percentile round trips through a same-zone endpoint?
\item And through the endpoint in another zone?
\item How much does the wrong zone add to each?
\item How does that compare with the step from peering to a same-zone endpoint?
\end{enumerate}

\medskip\noindent\textbf{Part II --- The cost.}
\begin{enumerate}[resume]
\item What does a million messages cost through each endpoint?
\item How much of the difference is cross-zone transfer?
\item What do ten billion messages a month cost through each?
\item What do the endpoint-hours cost in one zone and in three?
\end{enumerate}

\medskip\noindent\textbf{Part III --- The fix.}
\begin{enumerate}[resume]
\item How does the firm find out that it is in the wrong zone?
\item What are its options if the venue's endpoint has no interface in its zone?
\item What would a shared cluster \omterm{def:nw:trading-in-a-public-cloud:pg}{placement group} give it?
\item What does \cref{prop:nw:private-connectivity-to-crypto-venues:zone} say about moving the endpoint instead of the instances?
\end{enumerate}

\medskip\noindent\textbf{Part IV --- The verdict.}
\begin{enumerate}[resume]
\item State the \emph{named result}: the added latency and data-processing cost per million messages of a \omterm{def:nw:private-connectivity-to-crypto-venues:endpoint}{private endpoint} served from another zone
  against the same zone.
\item Which inputs are documentation, which measurement, which assumption?
\item Why is private connectivity's price not a reason to hesitate?
\item What should the venue publish to prevent the mistake?
\item How often should the firm re-check?
\item What does this problem have in common with chapter~9's cable coils?
\item What would change for a venue in a data centre?
\item In one sentence: what does a \omterm{def:nw:private-connectivity-to-crypto-venues:endpoint}{private endpoint} guarantee, and what does it not?
\end{enumerate}
\end{problem}

\section{Interview questions}

\begin{interviewq}[$\star$ \iqroles{developer}]\label{iq:nw:private-connectivity-to-crypto-venues:1}
What is the difference between VPC peering and PrivateLink?
\end{interviewq}

\begin{interviewq}[$\star\star$ \iqroles{developer}]\label{iq:nw:private-connectivity-to-crypto-venues:2}
How would you verify that your private connection to a venue does not cross an availability-zone boundary?
\end{interviewq}

\begin{interviewq}[$\star\star$ \iqroles{developer, trader}]\label{iq:nw:private-connectivity-to-crypto-venues:3}
A crypto venue offers four access tiers. How do you choose, and what do you ask for?
\end{interviewq}

\begin{interviewq}[$\star\star$ \iqroles{developer}]\label{iq:nw:private-connectivity-to-crypto-venues:4}
What adds latency in a private-endpoint path, and how would you measure each piece?
\end{interviewq}

\begin{interviewq}[$\star\star$ \iqroles{trader}]\label{iq:nw:private-connectivity-to-crypto-venues:5}
Is a shared \omterm{def:nw:trading-in-a-public-cloud:pg}{placement group} with the venue fair to other market participants?
\end{interviewq}

\begin{interviewq}[$\star\star\star$ \iqroles{developer}]\label{iq:nw:private-connectivity-to-crypto-venues:6}
Design the connectivity to a venue that runs in two zones of one region, for both speed and resilience.
\end{interviewq}
